\clearpage
\SourceLeaf{276}
\renewcommand{\BookRunningTitle}{Démonstration de quelques théorèmes}
\pdfbookmark[0]{Démonstration de quelques théorèmes concernants les périodes}{periodes-classes}
\begin{center}
\Large\bfseries DÉMONSTRATION DE QUELQUES THÉORÈMES CONCERNANTS\par\medskip
\large LES PÉRIODES DES CLASSES DES FORMES BINAIRES DU SECOND DEGRÉ.
\end{center}
\BookRule{.18\textwidth}

\pdfbookmark[1]{Théorème I}{periodes-theoreme-I}
\textsc{Théorème. I.}\quad\emph{Le nombre des classes (pr. pr.) d’un même déterminant, qui élevées à la dignité $P^{\mathrm{me}}$, $P$ étant ou un nombre premier ou la puissance d’un nombre premier $=p^\pi$, produisent la classe principale $K$, est égal ou à $1$ ou à une puissance de ce même nombre premier $p$.}

\emph{Démonstration.}\quad Soit $(\mathfrak Q)$ le groupe entier de toutes les classes en question et $n$ leur nombre. Puisque la classe principale $K$ est nécessairement contenue dans $(\mathfrak Q)$, le théorème est évident, si elle y est la seule. Mais s’il y en a d’autres, le nombre des classes contenues dans la période de chacune sera une puissance de $p$; soit une d’elles $A$, et supposons que sa période $(\mathfrak A)$ contienne $p^\alpha$ classes, qui seront toutes comprises dans $(\mathfrak Q)$. Or si les classes de cette période $(\mathfrak A)$ épuisent $(\mathfrak Q)$, on aura $p^\alpha=n$, et le théorème sera démontré; sinon, soit $B$ une classe quelconque de $(\mathfrak Q)$ non contenue dans $(\mathfrak A)$, et supposons que sa période soit développée jusqu’à ce qu’on y parvienne à une classe $bB$, qui soit en même temps parmi les classes de $(\mathfrak A)$, ce qui doit nécessairement arriver, parceque du moins la classe principale est commune à cette période et à $(\mathfrak A)$. Or supposant que $bB$ soit la \emph{première} classe dans la période de $B$ commune à $(\mathfrak A)$, ou $b$ le plus petit possible, je dis

$1^0$.\quad Que $b$ sera une puissance de $p$. Car il est évident qu’en faisant $b=p^\beta h$, $bB=iA$ et $hk\equiv1\pmod{p^\pi}$ (ce qui se pourra) on aura $kbB=p^\beta hkB=p^\beta B=ikA$,
\SourceLeaf{277}
c’est à dire que $p^\beta B$ sera aussi parmi les classes de $(\mathfrak A)$, d’où il s’ensuit que $h=1$ et $b=p^\beta$.

$2^0$.\quad Qu’en désignant les classes $K,B,2B\dotsfour(b-1)B$ par $(\mathfrak B)$, toutes les compositions d’une classe de $(\mathfrak A)$ avec une classe de $(\mathfrak B)$ donneront $p^{\alpha+\beta}$ classes différentes. Car en supposant $mA+nB=m'A+n'B$ et $n=n'$, on aura nécessairement $m=m'$; si $n>n'$, on aura $(n-n')B=(m'-m)A$, ce qui est impossible, si l’on n’a pas $n=n'$.

$3^0$.\quad Que ces $p^{\alpha+\beta}$ classes différentes seront comprises sous $(\mathfrak Q)$, ce qui est évident.

Or, si ces $p^{\alpha+\beta}$ classes épuisent $(\mathfrak Q)$, le théorème est démontré; sinon, on choisira nne autre classe de $(\mathfrak Q)$ non contenue parmi celles-là, savoir $C$; on continuera sa période jusqu’à ce qu’on y parvienne à une classe déjà comprise sous les classes composées de $(\mathfrak A)$ et $(\mathfrak B)$. Par un raisonnement semblable au précédent\textsuperscript{'} on démontrera, que l’exposant de cette classe doit être une puissance de $p$, $=p^\gamma$, et que la composition des $p^\gamma$ classes premières de la période de $C$ avec les $p^{\alpha+\beta}$ classes déjà trouvées donnera $p^{\alpha+\beta+\gamma}$ classes différentes toutes comprises dans $(\mathfrak Q)$. Si ces classes n’épuisent pas encore $(\mathfrak Q)$, on traitera de la même manière une quatrième classe $D$ etc., et il est évident que $(\mathfrak Q)$ étant formé d’un nombre fini de classes, ces opérations finiront aussi et qu’on aura $n$ égal à une puissance de $p$.\quad C. Q. F. D.

\par\bigskip
\pdfbookmark[1]{Théorème II}{periodes-theoreme-II}
\textsc{Théorème. II.}\quad\emph{Le nombre de toutes les classes du genre principal étant exprimé par $a^\alpha b^\beta c^\gamma$ etc., $a,b,c$, dénotant des nombres premiers différents, il y aura dans ce genre $a^\alpha,b^\beta,c^\gamma$ etc. classes, qui étant élevées à la dignité $a^\alpha,b^\beta,c^\gamma$ etc. resp. produisent la classe principale.}

\emph{Démonstration}\quad Soient $A,A',A''$ etc. toutes les classes qui élevées à la dignité $a^\alpha$ produisent $K$ et $(\mathfrak A)$ leur totalité; de même $B,B',B''$ etc. $(\mathfrak B)$, $C,C',C''$, $(\mathfrak C)$ etc. etc. Je dis que de la composition de toutes les classes de $(\mathfrak A)$ avec toutes les classes de $(\mathfrak B)$ avec toutes les classes de $(\mathfrak C)$ etc. il proviendra des classes différentes entre elles. Car si $A+B+C\dots=A'+B'+C'\dots$ etc., on aura, en faisant $A-A'=A''$, $B-B'=B''$ etc.,
\[
A''+B''+C''\text{ etc.}=K
\]
donc élevant à la dignité $b^\beta c^\gamma$ etc., $(b^\beta c^\gamma\dots)A''=K$, d’où il s’ensuit facilement
\SourceLeaf{278}
$A''=K$ et $A=A'$ et de la même manière on aura $B=B'$, $C=C'$ etc. Soit la totalité de ces classes $=(S)$. De plus il est clair que toutes ces classes seront du genre principal. Enfin il ne peut exister aucune classe dans le genre principal qui ne soit comprise sous $(S)$. Soit\ .\ .\ .
\BookRule{.34\textwidth}

\par\bigskip
\pdfbookmark[1]{Bemerkung}{periodes-bemerkung}
\begin{center}\textsc{Bemerkung.}\end{center}
\begingroup\small
Dieses im Jahre 1801 geschriebene Fragment bezieht sich auf Disq. Arithm. art. 306, \textsc{ix}. Das Wort \emph{dignité} wird hier in einem sonst nicht üblichen Sinne gebraucht.
\begin{flushright}\textsc{Stern}.\end{flushright}
\endgroup
\BookRule{.34\textwidth}

